\chapter{Literature Review}

\section{Classical Lambek calculus and completeness}
The classical literature provides an important baseline for this thesis. The original system of
\citet{Lambek1958} and \citet{Lambek1961} was quickly recognised as more than a linguistic
notation: it is a mathematically robust substructural logic with a rich proof theory and model
theory. Surveys and handbook treatments such as \citet{vanBenthem1988} and
\citet{Moortgat1997} show how the calculus sits within the wider landscape of categorial and
type-logical grammar.

For the specific question of completeness, the paper by \citet{Dosen1985} is especially relevant
because it directly establishes a completeness theorem for the Lambek calculus of syntactic
categories. Equally important is the relational-semantics perspective of
\citet{AndrekaMikulas1994}, who showed that relational models yield strong completeness results
for suitable versions of the calculus. From a language-theoretic viewpoint, \citet{Pentus1995}
proved completeness with respect to language models, while \citet{Pentus1993} and
\citet{Pentus1997} clarified the relation between Lambek grammars and context-free grammars.

These results matter for the present thesis in two ways. First, they demonstrate that
completeness questions for Lambek systems are deep but tractable. Second, they show that
there are already several mathematically natural semantic classes for which completeness holds.
This makes it reasonable to ask whether a similar story can be told for a probabilistic,
corpus-based semantics.


\section{Substructural logic and algebraic background}
The Lambek calculus belongs to the broader family of substructural logics, where structural
rules are controlled rather than assumed. This family includes relevance logics, linear logics, and
other resource-sensitive systems. Standard algebraic treatments organise these systems around
residuated structures, in which multiplication and division are linked by residuation principles
\citep{GalatosJipsenKowalskiOno2007,DosenSchroederHeister1993}.

This background is helpful because the semantics proposed in this thesis is explicitly designed to
mirror residuation: the multiplication clause aggregates compatible decompositions of a string,
while the two division clauses are defined by corpus-relative ratios. The result is not a standard
residuated algebra, but it is clearly inspired by the same adjoint pattern.


\section{Lambek calculus in linguistics and computational semantics}
Within linguistics, Lambek-style systems are valued because they combine lexical sensitivity,
word order, and compositionality in a transparent way. Textbook and monograph treatments
such as \citet{Morrill1994} and \citet{Moortgat1997} develop the linguistic applications in detail.
Steedman's work on CCG shows how categorial methods support wide-coverage syntactic and
semantic analysis \citep{Steedman2000}.

A further line of work connects categorial grammar to compositional models of meaning. While
this thesis is not primarily about vector semantics, the compositional viewpoint remains relevant:
if syntactic derivations guide semantic composition, then a probabilistic ranking over derivations
can also be interpreted as a ranking over candidate semantic analyses. For broader motivation on
composition-driven semantics, see \citet{CoeckeSadrzadehClark2010}.


\section{Probabilistic parsing: PCFGs, probabilistic CCGs, and stochastic Lambek grammars}
The most direct computational motivation for this thesis comes from probabilistic parsing.
Probabilistic context-free grammars are the standard starting point for ranking parses in the
presence of ambiguity \citep{ManningSchutze1999}. They are attractive because they combine a
clear formal semantics with trainable numerical parameters.

Probabilistic CCG takes similar ideas into a more lexicalised grammar formalism.
\citet{ClarkCurran2007} developed a wide-coverage statistical CCG parser using log-linear models,
showing that categorial methods can be scaled successfully when combined with probabilistic
ranking. For the Lambek setting more directly, \citet{BonfanteDeGroote2004} proposed a
stochastic point of view on Lambek categorial grammars, motivated precisely by the fact that
parsing in the Lambek calculus is inherently non-deterministic.

These works suggest that a probabilistic Lambek semantics is not an isolated curiosity. It sits at
a natural intersection of logical grammar and statistical parsing.


\section{Recent model-theoretic developments}
Recent work by Stepan Kuznetsov and collaborators shows that semantics and completeness for
Lambek-style systems remain active research topics. Of particular relevance is
\citet{Kuznetsov2023}, which revisits relational semantics for extensions of the Lambek calculus
and carefully distinguishes weak from strong completeness. Likewise,
\citet{KanovichKuznetsovScedrov2022} studies language models for extensions of the Lambek
calculus and highlights how adding connectives or constants can significantly affect completeness
behaviour.

For this thesis, the lesson is methodological: completeness is subtle even in well-studied,
non-probabilistic extensions. That makes it plausible that the completeness problem for a
corpus-relative probabilistic semantics will require new ideas rather than a straightforward
adaptation of existing proofs.


\section{Position of the present thesis}
The literature therefore leaves a clear space for the present project. On one side we have a
well-developed proof theory and completeness theory for classical Lambek calculi. On the other
side we have successful probabilistic grammar formalisms, together with some directly related
stochastic Lambek work. What is still missing is a careful, thesis-length account of a
corpus-based probabilistic semantics that is explicit about its formal choices, honest about its
limitations, and rigorous at least on soundness. That is the niche this thesis aims to fill.

